\documentclass[11pt]{article}

% Change "review" to "final" to generate the final (sometimes called camera-ready) version.
% Change to "preprint" to generate a non-anonymous version with page numbers.
\usepackage[]{acl}
\usepackage{times}
\usepackage{latexsym}

% For proper rendering and hyphenation of words containing Latin characters (including in bib files)
\usepackage[T1]{fontenc}
% For Vietnamese characters
% \usepackage[T5]{fontenc}
% See https://www.latex-project.org/help/documentation/encguide.pdf for other character sets

% This assumes your files are encoded as UTF8
\usepackage[utf8]{inputenc}
\usepackage{microtype}
\usepackage{inconsolata}
\usepackage{graphicx}
\usepackage{float}
\usepackage{amsmath}

% If the title and author information does not fit in the area allocated, uncomment the following
%
%\setlength\titlebox{<dim>}
%
% and set <dim> to something 5cm or larger.

\title{Group 10 Progress Report:\\Expected Goals (xG) Prediction in Soccer}


\author{Ali Zia, Andrew Habib, Third Author \\
  \texttt{\{ziam8,habiba21,macid3\}@mcmaster.ca} }

%\author{
%  \textbf{First Author\textsuperscript{1}},
%  \textbf{Second Author\textsuperscript{1,2}},
%  \textbf{Third T. Author\textsuperscript{1}},
%  \textbf{Fourth Author\textsuperscript{1}},
%\\
%  \textbf{Fifth Author\textsuperscript{1,2}},
%  \textbf{Sixth Author\textsuperscript{1}},
%  \textbf{Seventh Author\textsuperscript{1}},
%  \textbf{Eighth Author \textsuperscript{1,2,3,4}},
%\\
%  \textbf{Ninth Author\textsuperscript{1}},
%  \textbf{Tenth Author\textsuperscript{1}},
%  \textbf{Eleventh E. Author\textsuperscript{1,2,3,4,5}},
%  \textbf{Twelfth Author\textsuperscript{1}},
%\\
%  \textbf{Thirteenth Author\textsuperscript{3}},
%  \textbf{Fourteenth F. Author\textsuperscript{2,4}},
%  \textbf{Fifteenth Author\textsuperscript{1}},
%  \textbf{Sixteenth Author\textsuperscript{1}},
%\\
%  \textbf{Seventeenth S. Author\textsuperscript{4,5}},
%  \textbf{Eighteenth Author\textsuperscript{3,4}},
%  \textbf{Nineteenth N. Author\textsuperscript{2,5}},
%  \textbf{Twentieth Author\textsuperscript{1}}
%\\
%\\
%  \textsuperscript{1}Affiliation 1,
%  \textsuperscript{2}Affiliation 2,
%  \textsuperscript{3}Affiliation 3,
%  \textsuperscript{4}Affiliation 4,
%  \textsuperscript{5}Affiliation 5
%\\
%  \small{
%    \textbf{Correspondence:} \href{mailto:email@domain}{email@domain}
%  }
%}

\begin{document}
\maketitle
% \begin{abstract}
% \end{abstract}

\section{Introduction}

% Here, write a brief introduction to the problem you are solving. This can be adapted from your problem description and motivation from the original proposal. This should be around 0.25-0.5 pages.

Expected Goals (xG) is a frequently studied metric in the realm of modern soccer analytics utilized by various teams, sports analysts, and commentators to evaluate the quality of a scoring chance independent of the observed outcome. It predicts the probability of a shot in a soccer match resulting in a goal given the surrounding context used as features (examples include location, distance to goal, shooting angle, shot type, body part, technique, nearby opponents). This estimation can support analysis of team performance, player finishing, and defensive vulnerabilities while influencing future team strategies.

The xG prediction problem falls under binary classification where each shot is labelled as one of two classes using the single-label approach (goal or non-goal). Rather than purely focusing on the predicted class, the objective is to produce a value depicting the probability of scoring between 0-100 percent based on qualitative and quantitative features. Given a threshold such as 50 percent, one can predict if it is a goal or not.

The most significant challenge in this problem is that most of the dataset is labelled as no goal so a classifier predicting no goal can achieve a very high accuracy without truly evaluating chance quality. For instance, while two shots may have similar distances and angles, they can yield a different outcome due to the defender's position.

For this report, logistic regression is implemented along with binary cross-entropy through mini-batch gradient descent to relate shot characteristics to scoring probability. The performance of the best model in validation is compared to simple baselines Majority and Constant (0\% for every shot and predicts the average of goal labels in training data respectively) as well as the xG probabilities given by StatsBomb. The test set remains reserved for final evaluation on unseen data.

\section{Related Work}

Here, talk about the related work you encountered for your approach. Cite at least 5 references. Refer to item 2. No one has done exactly your task? Write about the most similar thing you can find. This should be around 0.25-0.5 pages.

\section{Dataset}
% You should write about your dataset here, following the guidelines regarding item 1. This section may be 0.5-1 pages. Depending on your specific dataset, you may want to include subsections for the preprocessing, annotation, etc.
\subsection{StatsBomb Open Data}
The dataset used in this project is \href{https://github.com/hudl/open-data}{StatsBomb Open Data} which is a free public collection of soccer event data hosted on GitHub. The dataset was collected via StatsBomb analysts going through match footage and tagging every action involving the ball (i.e. what happened, where it happened on the pitch and when it happened). Each match is stored as its own JSON file exported from the StatsBomb API. The data extraction process did not require web scraping or API access as the data is openly available. As a result, we cloned the repository directly and use the data under the StatsBomb Public Data User Agreement.
The current release of the dataset has 4,235 matches which is more than the 3,235 we identified in our proposal since the repository has grown since then. These matches come from 24 men's and women's competitions, including the top five European leagues, international tournaments and the FA Women's Super League.

\subsection{Labelling}
Each match file has around 3,500 events, with only roughly around 25 of them being shots. Every shot comes with an outcome recorded by StatsBomb, which can be Off Target, Blocked, Saved, Goal, Wayward, Post, Saved Off Target, or Saved to Post. We turned these outcomes into a binary label called \texttt{is\_goal}, which is 1 when the outcome is Goal and 0 for everything else. For the labelling process no manual annotation was done or needed as the labels come straight from what actually happened to each shot.

\subsection{Preprocessing}
Our preprocessing script \texttt{data\_pipeline.py} goes through every match file and pulls out all 108,281 shots into a single table, with one row per shot. We removed 363 penalty shoot-out attempts because they are not part of regular play (i.e. over time) and do not actually count towards the final score of the game. We also removed 7 shots that had a recorded location outside the $120\times80$ yard pitch. The resulting preprocessing left us with 107,911 shots in the final dataset.
Additionally, StatsBomb only includes true or false details such as \textit{first\_time} when they are true, so any missing value we set to 0. The 1,065 shots without a freeze frame (snapshot of where every player was standing when the shot was taken) are all penalties and a few other frames do not show the goalkeeper. These cases were filled in with fixed values instead of averages so that no information from the test set could carry over into training, and a \texttt{has\_freeze\_frame} flag was added to mark them. Any field that is only known after the ball is struck, like the outcome, where the shot ended up and StatsBomb's own xG value was left out of the model inputs since using it would give away the answer. Finally, the categorical details were one-hot encoded, giving 40 model inputs in total.

\subsection{Properties}
The final dataset has 107,911 shots. From the cumulative shot data 11,838 shots (11.0\%) are goals. Table~\ref{tab:dataset} shows that the chance of scoring drops quickly as the distance from goal increases and shots with no defender in the way are scored 2.6X as often as shots with at least one defender blocking the path.

\begin{table}[h]
  \centering
  \small
  \begin{tabular}{lrr}
    \hline
    \textbf{Subset} & \textbf{Shots} & \textbf{Goal rate} \\
    \hline
    All shots           & 107,911 & 11.0\% \\
    Open play           & 101,911 & 10.4\% \\
    Free kick (direct)  & 4,773   & 6.5\%  \\
    Penalty             & 1,194   & 76.5\% \\
    \hline
    Distance $\le$ 6 yd & 3,400   & 44.8\% \\
    6 to 12 yd          & 24,011  & 20.2\% \\
    12 to 18 yd         & 25,562  & 12.6\% \\
    18 to 24 yd         & 23,260  & 5.8\%  \\
    $>$ 24 yd           & 31,678  & 2.8\%  \\
    \hline
    No defender in path & 51,885  & 16.1\% \\
    $\ge$ 1 defender    & 56,026  & 6.2\%  \\
    \hline
  \end{tabular}
  \caption{\label{tab:dataset}
    Goal rate by shot type, distance to goal, and whether an outfield defender is standing between the shooter and the goal.}
\end{table}

\section{Features}

Describe any features you used for your model, or how your data was input to your model. Are you doing feature engineering or feature selection? Are you learning embeddings? Is it all part of one neural network? Refer to item 2. This may range from 0.25 pages to 0.5 pages.

\section{Implementation}

% Describe your model and implementation here. Refer to item 4. This may take around a page.

The implementation stage involved manually implementing a logistic regression model using NumPy and the techniques taught in SFWRENG 4AL3 Tutorial 3 to predict the probability that a soccer shot results in a goal. Each shot observation is represented by the 40 features described in the Features section. The target label is binary, representing goal (\(y=1\)) or no goal (\(y=0\)). Rather than simply returning a class label, the model returns a probability between 0 and 1, representing its expected-goals prediction.

Before training begins, features are standardized using the training-set means and standard deviations, with operations applied element-wise. The same training-set values are used to scale validation inputs. Any zero standard deviation is replaced with 1 to avoid division by zero:
\[
\mathbf{x}_{\mathrm{scaled}} = (\mathbf{x} - \boldsymbol{\mu}_{\mathrm{train}})/\boldsymbol{\sigma}_{\mathrm{train}}.
\]

For an input vector \(\mathbf{x}\), the linear model produces a score, where the sigmoid function maps the value to a probability ranging between 0 and 1:
\[
z = \mathbf{w} \cdot \mathbf{x}_{\mathrm{scaled}} + b,
\]
\[
\hat{y}=P(y=1\mid\mathbf{x})=\sigma(z)=\frac{1}{1+e^{-z}},
\]

\(\mathbf{x}_{\mathrm{scaled}}\) contains the standardized features, \(\mathbf{w}\) is the vector of learned weights, and \(b\) is the learned bias. \(\hat{y}\) represents the predicted probability of a goal, while \(y\) is the actual outcome: 1 meaning goal and 0 meaning non-goal.

The model is trained by minimizing average binary cross-entropy loss. This loss function is useful beyond just evaluating whether a model correctly predicts goal or no goal since it rewards assigning higher probabilities to the correct class. The same implementation is used on both training and validation losses. Training loss measures how well the model fits the training data and influences weights and biases during gradient descent steps while validation loss is used to select the best performing model given a set of hyperparameters. The stable CE implementation was used to avoid overflow errors:
\[
\begin{aligned}
L_{\mathrm{CE}}
&= -\frac{1}{N}\sum_{i=1}^{N}
\left[y_i\log(\hat{y}_i)+(1-y_i)\log(1-\hat{y}_i)\right] \\
&= \frac{1}{N}\sum_{i=1}^{N}
\left[\log(1+e^{z_i})-y_i z_i\right].
\end{aligned}
\]

The model is optimized using mini-batch stochastic gradient descent. At the beginning of each epoch (one pass through the training set), the training observations are shuffled and split into mini-batches. Every training data point is used once per epoch. For each mini-batch, predictions are calculated using the current parameters, and the average cross-entropy gradients are computed following Tutorial 3:
\[
\begin{aligned}
\frac{\partial L_{\mathrm{CE}}(\hat{y}, y)}{\partial w_j}
&= \frac{1}{n}\sum_{i=1}^{n}
(\hat{y}^{(i)}-y^{(i)})x_j^{(i)}, \\
\frac{\partial L_{\mathrm{CE}}(\hat{y}, y)}{\partial b}
&= \frac{1}{n}\sum_{i=1}^{n}
(\hat{y}^{(i)}-y^{(i)}).
\end{aligned}
\]
\(n\) is the number of rows in the current mini-batch, and \(x_j^{(i)}\) is the standardized value of feature \(j\) for data point \(i\). Since the model has 40 features, the tutorial's weight gradient is calculated for each weight \(w_j\). The weights and bias are updated after each mini-batch:
\[
\begin{aligned}
w_j &\leftarrow w_j-\eta\frac{\partial L_{\mathrm{CE}}(\hat{y}, y)}{\partial w_j}, \\
b &\leftarrow b-\eta\frac{\partial L_{\mathrm{CE}}(\hat{y}, y)}{\partial b},
\end{aligned}
\]
where \(\eta\) is the learning rate. The next mini-batch uses these updated parameters. Training cross-entropy is recorded after each epoch to monitor learning. Validation loss is used only to compare different sets of hyperparameters.

A random search with 20 trials over 100 epochs. Learning rates were sampled uniformly in logarithmic space between \(0.0001\) and \(0.1\), and various batch sizes were chosen \(\{32,64,128,256,512,1024\}\). A seed of 42 was used for consistency. The set of hyperparameters that yielded the lowest validation cross-entropy was selected, giving a learning rate of approximately \(0.011198\) and a mini-batch size of 32.

\begingroup
\setlength{\intextsep}{4pt}
\begin{figure}[H]
\centering
\setlength{\abovecaptionskip}{3pt}
\setlength{\belowcaptionskip}{0pt}
\includegraphics[width=\columnwidth]{../progress-report/artifacts/training_loss.png}
\caption{Training cross-entropy over 100 epochs for the selected logistic regression model.}
\label{fig:training-loss}
\end{figure}
\endgroup

The Logistic Regression model was compared to majority vote which always predicts no goal, constant-probability baseline which predicts based on the full training set's average of binary labels, and the model produced by Statsbomb which uses gradient-boosted decision trees. Note: Statsbomb's validation loss comes from the same CE loss equation where \(\hat{y}_i\) is the Statsbomb xG probability extracted from the dataset, while \(y_i\) is the actual outcome both coming from the validation split.

\begin{table}[H]
\centering
\resizebox{\columnwidth}{!}{%
\begin{tabular}{lrr}
\hline
Model/reference & Training CE & Validation CE \\
\hline
Logistic regression & 0.269440 & 0.266354 \\
Majority vote & 3.025423 & 3.016689 \\
Constant probability & 0.345455 & 0.344793 \\
StatsBomb xG & --- & 0.264980 \\
\hline
\end{tabular}
}
\caption{Training and validation cross-entropy losses for the logistic regression, majority-class baseline, constant-probability baseline, and StatsBomb xG.}
\label{tab:validation-loss}
\end{table}

The results show that the model our team developed using Logistic Regression yields lower losses than the Majority vote and Constant probability baseline models by 2.750335 and 0.078439 respectively. The model also performs almost as well as Statsbomb's optimized model with a difference of about -0.001374.

\section{Results and Evaluation}

How are you evaluating your model? What results do you have so far? What are your baselines? Refer to item 5. This may take around 0.5 pages.

\section{Feedback and Plans}

Write about your plans for the remainder of the project. This should include a discussion of the feedback you received from your TA, and how you plan to improve your approach. Reflect on your implementation and areas for improvement. Refer to item 6. This may be around 0.5 pages.

\section{Template Notes}

You can remove this section or comment it out, as it only contains instructions for how to use this template. You may use subsections in your document as you find appropriate.

\subsection{Tables and figures}

See Table~\ref{citation-guide} for an example of a table and its caption.


\begin{figure*}[t]
  \includegraphics[width=0.48\linewidth]{example-image-a} \hfill
  \includegraphics[width=0.48\linewidth]{example-image-b}
  \caption {A minimal working example to demonstrate how to place
    two images side-by-side.}
\end{figure*}


\subsection{Citations}

\begin{table*}
  \centering
  \begin{tabular}{lll}
    \hline
    \textbf{Output}           & \textbf{natbib command} & \textbf{ACL only command} \\
    \hline
    \citep{Gusfield:97}       & \verb|\citep|           &                           \\
    \citealp{Gusfield:97}     & \verb|\citealp|         &                           \\
    \citet{Gusfield:97}       & \verb|\citet|           &                           \\
    \citeyearpar{Gusfield:97} & \verb|\citeyearpar|     &                           \\
    \citeposs{Gusfield:97}    &                         & \verb|\citeposs|          \\
    \hline
  \end{tabular}
  \caption{\label{citation-guide}
    Citation commands supported by the style file.
  }
\end{table*}

Table~\ref{citation-guide} shows the syntax supported by the style files.
We encourage you to use the natbib styles.
You can use the command \verb|\citet| (cite in text) to get ``author (year)'' citations, like this citation to a paper by \citet{Gusfield:97}.
You can use the command \verb|\citep| (cite in parentheses) to get ``(author, year)'' citations \citep{Gusfield:97}.
You can use the command \verb|\citealp| (alternative cite without parentheses) to get ``author, year'' citations, which is useful for using citations within parentheses (e.g. \citealp{Gusfield:97}).

\subsection{References}

\nocite{Ando2005,andrew2007scalable,rasooli-tetrault-2015}

Many websites where you can find academic papers also allow you to export a bib file for citation or bib formatted entry. Copy this into the \texttt{custom.bib} and you will be able to cite the paper in the \LaTeX{}. You can remove the example entries.

\subsection{Equations}

An example equation is shown below:
\begin{equation}
  \label{eq:example}
  A = \pi r^2
\end{equation}

Labels for equation numbers, sections, subsections, figures and tables
are all defined with the \verb|\label{label}| command and cross references
to them are made with the \verb|\ref{label}| command.
This an example cross-reference to Equation~\ref{eq:example}. You can also write equations inline, like this: $A=\pi r^2$.


% \section*{Limitations}

\section*{Team Contributions}

\textbf{Andrew Habib (habiba21):} Completed the Introduction and Implementation sections of the report and implemented the model and training pipeline while comparing to majority, constant, and statsbomb baselines.

\textbf{Ali Zia (ziam8):} Completed the Dataset and Results and Evaluations sections of the report and implemented the data preprocessing code.

% Bibliography entries for the entire Anthology, followed by custom entries
%\bibliography{custom,anthology-overleaf-1,anthology-overleaf-2}

% Custom bibliography entries only
\bibliography{custom}

% \appendix

% \section{Example Appendix}
% \label{sec:appendix}

% This is an appendix.

\end{document}
